\chapter{Strong IH: the complete local corner}
\label{component:pf3e0ca8bf054}
\noindent\textbf{Component }\path{pf3e0ca8bf054}\par
{\small\noindent Source: \path{R13/baseline/baseline/advances/ih/strong/ih_local_corner.tex}\par}
{\small\noindent SHA256: \path{f3e0ca8bf0547150036dcd1c624a329da3804b1a781c24835343af7acd1ba115}\par}
{\small\noindent Inventory use: SL1--SL5 strong-IH full scalar/current/normal construction on its own source bounds\par}
\medskip
\section{Strong IH: the complete local corner}
\label{pf3e0ca8bf054:sec:strong-ih-local}
Assume the strong IH source entry of \S\ref{pf5726218aec8:sec:strong-angle-entry}:
$z$ is infinitesimal, $y$ unbounded, $B=H+L$ a bounded positive unit,
$0<H<1$, $L>0$, and $C,\alpha$ strong. The source parameter
$\rho=B/y$ is infinitesimal, $s=r/\rho$ bounded, and $c_s=(1+L)/B$
is a positive unit. Testing $u,w=1/v$ are infinitesimal or zero.
We use the exact graph, five roots and both full normal identities
(IC3)--(IC12) of the preserved corner calculation; none of those
algebraic identities assumes $C$ weak.

The source quantities obey (GE3), and direct rational simplification
gives additionally
\[
 \frac{k}{B\rho}\longrightarrow
 \frac{2(B^2s+BHs+BH-Bs-H^2+H)}{B^2}.
\tag{SL1}
\]
No opening denominator occurs in this normalized expression.

Set $T_\rho=\rho T$ and
$\eta=\sqrt{H(1-H)/B}$, so $Y=HwZ+\eta E_Y$.
The full Gram's entire explicit dependence on $C$ is exactly
\[
 4C(U_1^2+U_2^2),\quad
 U_1=uW-Y-LwZ=\sqrt B R_0-\eta E_Y,
\]
\[
 U_2=w(uW-Y)-(L/y)Z
       =R_2-\sqrt{H/B}E-w\eta E_Y.
\tag{SL2}
\]
The scalar and normal identities verifying this assertion use
\[
 \sqrt H R_2-\sqrt{BH}\,wR_0-\sqrt B E=0.
\tag{SL3}
\]
Choose a positive real cap $c_*\le C$ and retain
$4(C-c_*)(U_1^2+U_2^2)$. In the remaining exact Gram replace explicit
$C$ by $c_*$, leaving $k$ and all other original source quantities
unchanged. The distinction matters: this is an exact square
subtraction from the complete Gram, not a perturbation of the source.
Add $\lambda$ times the exact zero (WIE2), where $\lambda>0$ is a
real number to be chosen. All resulting coefficients are bounded,
using (GE3), (SL1), and the opening-normalized roots.

We display the entire residue. The roots $R_1$ and $E_X$ have isolated
positive pivots $2c_s$ and $1$. The $(R_0,E_Y)$ block is
\[
 \begin{pmatrix}
 4Bc_*&-4c_*\sqrt{H(1-H)}\\
 -4c_*\sqrt{H(1-H)}&
 (4c_*+4A+4\alpha+2)H(1-H)/B
 \end{pmatrix}.
\tag{SL4}
\]
Here and below all expressions in a residue matrix mean their bounded
residues; products such as $AH$ are evaluated by (GE3), even when
$A$ itself is unbounded. Its determinant is exactly $8c_*(1+L)>0$.
Put $a_*=2(2c_*+1)$. The $(R_2,E)$ block is
\[
 \begin{pmatrix}
 a_*&-a_*\sqrt{H/B}\\
 -a_*\sqrt{H/B}&H(4A+a_*)/B
 \end{pmatrix},
\tag{SL5}
\]
with determinant $2a_*c_s>0$. The only remaining coupling is between
$R_2$ and $T_\rho$, whose Gram entry is
$c_s(H/B-s)$. The $T_\rho$ diagonal is $\lambda+2c_s s$.
Choose $\lambda$ larger than the finite Schur-complement threshold
of the positive block (SL5). This produces a strictly positive
residue matrix. Since this choice is made after taking a real
residue, it is a real scalar and introduces no mixed inverse.

We spell out why no normal term is hidden in this block description.
In (IC11), $f_X/\rho$ is bounded by (GE3) and every scalar-root
multiplier contains $u$ or $w$. In (IC12), the only nonzero limiting
term is the explicit $C$ contribution to $R_0E_Y$, already included
in (SL2)--(SL4). After the cap, all other normal multipliers contain
$z$, $u$, or $w$. The unbounded factors $A$ and $\alpha$ always occur
as $AH$ and $\alpha(1-H)$ after normalization, hence remain bounded
also at $H_0=0,1$. The normal $X$ coefficient divided by $N=4n+2$
has residue one, and its coupling to $E_Y$ vanishes. The remaining
$Y$ coefficient is exactly the lower-right entry of (SL4).
These facts follow directly from the displayed exact normal
factorizations, not from scalar-graph positivity.

The bounded positive-residue Gram lemma now gives tensor SOS for the
complete original parent, after restoring the retained positive
$C-c_*$ squares. This holds for all states and all independent local
tests, for both opening endpoints, for arbitrary $n\ge0$, and at
literal $u=0$ or $w=0$.
